% bismi ALLAH ARRAHMN ARRAHIM
























































































\documentclass[11pt]{article}

\usepackage[a4paper,margin=2.6cm]{geometry}
\usepackage{amsmath,amssymb,amsthm,mathtools,bm,booktabs,enumitem}
\usepackage[hidelinks]{hyperref}

\newtheorem{definition}{Définition}
\newtheorem{proposition}{Proposition}
\newtheorem{remark}{Remarque}

\providecommand{\Ppch}{\mathcal P^{\mathrm{PCH}}}
\providecommand{\Tpch}{\mathcal T^{\mathrm{PCH}}_{\eta}}
\providecommand{\Gain}{\mathsf{Gain}}

\title{Spécification rigoureuse d'un LLM purement PCH sans attention standard}
\author{}
\date{}

\begin{document}
\maketitle

\section{Objet et principe directeur}

On fixe ici une architecture de \emph{LLM purement PCH} : aucun bloc d'attention standard, aucun softmax sur des produits requêtes-clés, aucun mécanisme de type
\[
\mathrm{softmax}\!\left(\frac{QK^\top}{\sqrt d}\right)V.
\]
Toute la propagation d'information se fait par des blocs proximaux PCH, des résumés globaux non attentionnels, des MLP résiduels, et un schéma de diffusion en espace latent.

Le point de départ est le noyau PCH exact utilisé dans les expériences précédentes :
\[
\Ppch(w;c)
=
\arg\min_u
\left\{
\frac12\|u-w\|^2+\Phi_{\mathrm{PCH},c}(u)
\right\},
\]
avec
\[
\Phi_{\mathrm{PCH},c}(u)
=
\lambda_{\mathrm{ctx}}\operatorname{dist}(u,\mathcal M_c)^2
+
\lambda_{\mathrm{ch}}\Omega_{\mathrm{ch}}(u)
+
\lambda_{\mathrm{reg}}\|u\|^2.
\]
Dans la version expérimentale de référence, on instancie explicitement
\[
\Omega_{\mathrm{ch}}(u)=\|u\|^2,
\qquad
\lambda_{\mathrm{ridge}}=\lambda_{\mathrm{ch}}+\lambda_{\mathrm{reg}}.
\]
Le noyau est donc un prox exact d'une énergie convexe, résolu par programmation quadratique convexe.

\section{Notation séquentielle}

Soit une séquence de tokens
\[
x_{1:n}=(x_1,\dots,x_n).
\]
On note :
\[
e_i = E_{\theta}(x_i),\qquad
h_i^{(0)} = e_i + p_i,
\]
où \(E_\theta\) est l'embedding de tokens appris et \(p_i\) un encodage positionnel déterministe.

On regroupe les états dans une matrice :
\[
H^{(\ell)} = (h_1^{(\ell)},\dots,h_n^{(\ell)})^\top \in \mathbb R^{n\times d}.
\]

\section{Architecture générale du LLM PCH}

L'architecture est une pile de blocs sans attention standard :
\[
H^{(0)}
\;\longrightarrow\;
\mathsf{PCHBlock}_1
\;\longrightarrow\;
\mathsf{PCHBlock}_2
\;\longrightarrow\;\cdots\;\longrightarrow\;
\mathsf{PCHBlock}_L
\;\longrightarrow\;
\mathsf{Decoder}.
\]

Chaque bloc \(\mathsf{PCHBlock}_\ell\) comporte trois sous-étapes :
\begin{enumerate}[leftmargin=2em]
    \item un routage de contexte non attentionnel ;
    \item un prox PCH exact ;
    \item un MLP résiduel de mélange local.
\end{enumerate}

Il n'existe nulle part de mécanisme de type requête-clé-valeur.

\subsection{Réservoir de contextes}

À chaque couche \(\ell\), on maintient un réservoir de contextes
\[
\mathcal B_\ell = \{\mathcal B^{(\ell)}_c\}_{c=1}^{C_\ell},
\]
où \(\mathcal B^{(\ell)}_c\) est un ensemble de calibrations pour le contexte \(c\).

On définit le contexte admissible par l'enveloppe convexe
\[
\mathcal M^{(\ell)}_c
=
\operatorname{conv}\{\psi_\ell(v): v\in \mathcal B^{(\ell)}_c\},
\]
où \(\psi_\ell\) est une projection linéaire ou un petit réseau linéaire appris de faible rang, utilisé uniquement pour placer les calibrations dans l'espace latent de la couche \(\ell\).

Dans l'implémentation pure, le contexte n'est pas inféré par attention, mais par sélection géométrique :
\[
c_i^{(\ell)}
=
\arg\min_{c\in\{1,\dots,C_\ell\}}
\operatorname{dist}\!\left(h_i^{(\ell)}, \mathcal M^{(\ell)}_c\right)^2.
\]
C'est le routeur par distance de projection, qui est le routeur géométrique le plus direct et le plus fidèle au noyau testé expérimentalement.

\subsection{Bloc PCH d'une couche}

Pour chaque token \(i\) et chaque couche \(\ell\), on applique :
\[
\tilde h_i^{(\ell)}
=
\Ppch_\ell\!\left(h_i^{(\ell)}; c_i^{(\ell)}\right)
=
\arg\min_u
\left\{
\frac12\|u-h_i^{(\ell)}\|^2
+
\Phi_{\mathrm{PCH},c_i^{(\ell)}}^{(\ell)}(u)
\right\}.
\]

Ici
\[
\Phi_{\mathrm{PCH},c}^{(\ell)}(u)
=
\lambda_{\mathrm{ctx}}^{(\ell)}
\operatorname{dist}(u,\mathcal M_c^{(\ell)})^2
+
\lambda_{\mathrm{ch}}^{(\ell)}\Omega_{\mathrm{ch}}^{(\ell)}(u)
+
\lambda_{\mathrm{reg}}^{(\ell)}\|u\|^2.
\]

La sortie de la couche est ensuite donnée par un résiduel simple :
\[
h_i^{(\ell+1)}
=
\tilde h_i^{(\ell)}
+
\mathrm{MLP}_\ell\!\left(\tilde h_i^{(\ell)}\right).
\]
On peut ajouter une normalisation de couche si nécessaire, mais aucune attention n'est introduite.

\subsection{Interaction entre tokens sans attention}

Pour obtenir un couplage séquentiel sans attention, on utilise deux canaux seulement :

\begin{enumerate}[leftmargin=2em]
    \item un canal \emph{global} via un résumé
    \[
    s^{(\ell)}=\frac1n\sum_{i=1}^n h_i^{(\ell)},
    \]
    qui conditionne le routage de contexte ;
    \item un canal \emph{géométrique} via les hulls \(\mathcal M_c^{(\ell)}\), qui sont construits à partir de banques de calibration et de l'état latent courant.
\end{enumerate}

Ainsi, l'information circule par contraction géométrique vers des ensembles admissibles, et non par pondérations attentionnelles pair-à-pair.

\section{Version diffusion : génération non autoregressive}

Pour éviter le coût d'un décodage token-par-token, le modèle est formulé comme un \emph{modèle de diffusion en espace latent}. La génération et l'entraînement portent alors sur des séquences entières, en parallèle.

\subsection{Processus direct}

Soit \(Z_0 \in \mathbb R^{n\times d}\) le latent propre associé à une séquence de texte.
On définit un processus de noising classique :
\[
Z_t
=
\sqrt{\bar\alpha_t}\,Z_0
+
\sqrt{1-\bar\alpha_t}\,\varepsilon,
\qquad
\varepsilon\sim\mathcal N(0,I),
\]
où \(\bar\alpha_t\) suit un calendrier de diffusion fixé.

\subsection{Processus inverse PCH}

Le débruitage inverse n'utilise ni attention ni score-network à attention.
Il est porté par un débruiteur latent \(\mu_\theta\), puis corrigé par un pas PCH exact.

Le cœur du pas inverse est :
\[
Z_{t-1}
=
\arg\min_Z
\left\{
\frac1{2\sigma_t^2}\|Z-\mu_\theta(Z_t,t)\|_F^2
+
\Phi_{\mathrm{PCH},t}(Z)
\right\},
\]
où \(\Phi_{\mathrm{PCH},t}\) agit ligne par ligne sur les états latents des tokens, avec des contextes définis par la banque \(\mathcal B_t\).

Autrement dit, la diffusion inverse est un \emph{pas proximal corrigé} :
\[
Z_{t-1}
=
\mathsf{PCHReverse}(Z_t,t).
\]
Cette formulation est entièrement compatible avec le noyau PCH exact : le terme proximal est une vraie minimisation convexe, pas une approximation stochastique cachée.

\section{Architecture complète d'un LLM PCH pur}

La version complète est la suivante.

\subsection{Encodage}

\begin{enumerate}[leftmargin=2em]
    \item Tokenisation.
    \item Embedding appris \(E_\theta\).
    \item Encodage positionnel.
    \item Passage dans une pile de blocs PCH.
\end{enumerate}

\subsection{Bloc élémentaire}

Pour chaque couche \(\ell\), on exécute :
\[
H^{(\ell)}
\;\mapsto\;
C^{(\ell)}
\;\mapsto\;
\tilde H^{(\ell)}
\;\mapsto\;
H^{(\ell+1)}.
\]

Plus précisément :
\[
c_i^{(\ell)}
=
\arg\min_c \operatorname{dist}(h_i^{(\ell)}, \mathcal M_c^{(\ell)})^2,
\]
\[
\tilde h_i^{(\ell)}
=
\Ppch_\ell(h_i^{(\ell)};c_i^{(\ell)}),
\]
\[
h_i^{(\ell+1)}
=
\tilde h_i^{(\ell)}+\mathrm{MLP}_\ell(\tilde h_i^{(\ell)}).
\]

\subsection{Sortie}

La tête de sortie est un décodeur linéaire ou MLP :
\[
\mathrm{logits}_i = W_{\mathrm{out}} h_i^{(L)} + b_{\mathrm{out}}.
\]
Le texte est produit soit :
\begin{itemize}[leftmargin=2em]
    \item par débruitage diffusion puis décodage parallèle ;
    \item soit par raffinements successifs sur toute la séquence.
\end{itemize}

Il n'y a pas de génération autoregressive basée sur l'attention.

\section{Objectif d'entraînement}

L'entraînement combine un terme de diffusion et un terme de reconstruction de texte.

\subsection{Perte diffusion}

On minimise :
\[
\mathcal L_{\mathrm{diff}}
=
\mathbb E_{x,t,\varepsilon}
\left[
\left\|
Z_0-\widehat Z_0(Z_t,t)
\right\|_F^2
\right],
\]
où \(\widehat Z_0\) est la prédiction du latent propre par le réseau PCH.

\subsection{Perte textuelle}

On ajoute une perte de décodage token :
\[
\mathcal L_{\mathrm{tok}}
=
-\sum_{i=1}^n \log p_\theta(x_i \mid H^{(L)}).
\]

\subsection{Pertes de cohérence PCH}

On ajoute des régularisations compatibles avec le noyau :
\[
\mathcal L_{\mathrm{PCH}}
=
\sum_{\ell=1}^L\sum_{i=1}^n
J_{\mathrm{PCH}}^{(\ell)}\!\left(\tilde h_i^{(\ell)}; h_i^{(\ell)}, c_i^{(\ell)}\right),
\]
ainsi qu'un terme de stabilité du routage si besoin :
\[
\mathcal L_{\mathrm{route}}
=
\sum_{\ell,i}
\operatorname{dist}\!\left(h_i^{(\ell)}, \mathcal M_{c_i^{(\ell)}}^{(\ell)}\right)^2.
\]

La perte totale est
\[
\mathcal L
=
\mathcal L_{\mathrm{diff}}
+
\beta\,\mathcal L_{\mathrm{tok}}
+
\gamma\,\mathcal L_{\mathrm{PCH}}
+
\delta\,\mathcal L_{\mathrm{route}}.
\]

\section{Protocoles d'implémentation}

\subsection{Bootstrap}

Au démarrage, on initialise :
\begin{itemize}[leftmargin=2em]
    \item l'embedding \(E_\theta\),
    \item les MLP de couche,
    \item la banque de contextes \(\mathcal B_\ell\) par calibrations initiales,
    \item les projections \(\psi_\ell\).
\end{itemize}

Les contextes peuvent être construits par un pré-clustering géométrique des états latents initiaux.

\subsection{Boucle d'entraînement}

Pour chaque mini-batch :
\begin{enumerate}[leftmargin=2em]
    \item encoder la séquence en latent propre \(Z_0\) ;
    \item échantillonner un temps \(t\) ;
    \item corroder \(Z_0\) en \(Z_t\) ;
    \item appliquer la pile de blocs PCH sur \(Z_t\) ;
    \item reconstruire \(\widehat Z_0\) et les logits textuels ;
    \item calculer \(\mathcal L\) ;
    \item mettre à jour les paramètres par rétropropagation à travers les solveurs convexes différentiables ;
    \item rafraîchir périodiquement les contextes \(\mathcal B_\ell\) à partir des activations récentes.
\end{enumerate}

\subsection{Rafraîchissement des contextes}

Les ensembles \(\mathcal B_\ell\) sont mis à jour par buffer glissant, de manière périodique et contrôlée.
On conserve pour chaque couche des calibrations représentatives, puis on recalcule les hulls :
\[
\mathcal M_c^{(\ell)}
=
\operatorname{conv}\{\psi_\ell(v): v\in\mathcal B_c^{(\ell)}\}.
\]
Ce mécanisme garde le noyau PCH fidèle à sa forme exacte tout en lui permettant d'opérer sur des représentations apprises.

\subsection{Inférence}

À l'inférence :
\begin{enumerate}[leftmargin=2em]
    \item on part d'un bruit latent ou d'une représentation partielle ;
    \item on itère le processus inverse PCH ;
    \item on décode en sortie ;
    \item on peut appliquer un second passage de purification si nécessaire.
\end{enumerate}

\section{Protocole de validation}

La validation d'un LLM PCH pur doit vérifier quatre familles de propriétés :

\begin{enumerate}[leftmargin=2em]
    \item \textbf{stabilité proximale} :
    décroissance d'énergie et résidus de point fixe faibles ;
    \item \textbf{stabilité de diffusion} :
    reconstruction cohérente du latent propre ;
    \item \textbf{cohérence séquentielle} :
    amélioration de perplexité et de qualité de génération ;
    \item \textbf{cohérence géométrique} :
    maintien d'une contraction locale et réduction de distance aux hulls admissibles.
\end{enumerate}

\section{Complexité et mise en oeuvre}

Le coût dominant vient des sous-problèmes convexes PCH.
Pour une séquence de longueur \(n\), une couche et un pas diffusion, on résout \(n\) prox exacts, chacun conditionné par un contexte.

Le schéma reste donc :
\[
\text{séquence} \to \text{projections convexes} \to \text{MLP} \to \text{diffusion inverse}.
\]
Il est massivement parallélisable sur les tokens et sur les couches, à condition de réutiliser des banks de contextes partagées et des solveurs warm-startés.

\section{Garantie de fidélité au noyau PCH}

Cette architecture respecte le noyau PCH exact dans le sens suivant :
\begin{itemize}[leftmargin=2em]
    \item la purification contextuelle est toujours définie comme un vrai minimiseur convexe ;
    \item le contexte est un hull convexe réel ;
    \item la monotonie énergétique peut être contrôlée couche par couche ;
    \item la contraction locale reste mesurable par linéarisation ;
    \item le routage est séparé du noyau et ne remplace pas le prox.
\end{itemize}

Le modèle est donc un \emph{LLM à noyau proximal contextuel}, non un Transformer déguisé.

\section{Résumé opérationnel}

Le pipeline complet est :
\[
x_{1:n}
\;\to\;
E_\theta
\;\to\;
H^{(0)}
\;\to\;
\mathsf{PCHBlock}_1
\;\to\;\cdots\;\to\;
\mathsf{PCHBlock}_L
\;\to\;
\mathsf{DiffusionHead}
\;\to\;
\text{texte}.
\]

Chaque bloc PCH fait :
\[
\text{routage géométrique}
\;\to\;
\text{prox exact}
\;\to\;
\text{résiduel MLP}.
\]

Aucune attention standard n'intervient à aucun endroit.



%%%%%%%%%%%%%%%%%%%%%%%%%%%%%%%%%%%%%%%%%%%%%%%%%%%%%%%%%%%%%%%%%%%%%%%%%%%%%%%%%%%%%%%%%%%%%%%pipeline de bout en bout 

\section{Pipeline complet d'un LLM PCH pur : de l'entraînement à l'usage chat}

\subsection{Vue d'ensemble du pipeline}

On considère un LLM purement PCH dont le cycle de vie complet est le suivant :
\[
\text{données textuelles}
\;\to\;
\text{tokenisation}
\;\to\;
\text{encodage latent}
\;\to\;
\text{entraînement conjoint PCH + diffusion}
\;\to\;
\text{sauvegarde de checkpoint}
\;\to\;
\text{chargement du modèle entraîné}
\;\to\;
\text{usage chat prompt/réponse}.
\]

Le point essentiel est que le modèle ne comporte à aucun endroit de mécanisme d'attention standard. Toute propagation d'information repose sur :
\begin{itemize}[leftmargin=2em]
    \item des encodages latents appris ;
    \item des résumés globaux sans attention, par pooling ou agrégation géométrique ;
    \item des blocs PCH exacts ;
    \item une dynamique de diffusion en espace latent ;
    \item un décodeur final non attentionnel.
\end{itemize}

\subsection{Format des données d'entraînement}

Les données sont organisées sous forme de paires \((p, r)\), où :
\[
p = \text{prompt / contexte dialogue},
\qquad
r = \text{réponse cible}.
\]

Dans le cas d'un pré-entraînement sur corpus brut, on peut aussi utiliser des segments textuels longs :
\[
x = (x_1,\dots,x_n),
\]
puis former des tâches de reconstruction, de débruitage ou de continuation structurée.

Pour un usage conversationnel, on définit une séquence complète :
\[
s = (\texttt{system}, \texttt{user}_1, \texttt{assistant}_1, \dots, \texttt{user}_k, \texttt{assistant}_k),
\]
puis on isole :
\[
p = \text{historique de dialogue + message courant},
\qquad
r = \text{réponse à générer}.
\]

\subsection{Tokenisation et encodage}

Soit un tokenizer déterministe \(\tau\). On obtient une séquence de tokens :
\[
\tau(s) = (t_1,\dots,t_n).
\]

Chaque token est transformé en embedding :
\[
e_i = E_\theta(t_i),
\]
où \(E_\theta\) est l'embedding appris.

On ajoute un encodage positionnel fixe ou appris :
\[
h_i^{(0)} = e_i + p_i.
\]

La matrice initiale est donc :
\[
H^{(0)} = (h_1^{(0)},\dots,h_n^{(0)})^\top \in \mathbb R^{n\times d}.
\]

Aucune attention n'intervient dans cette étape.

\subsection{Organisation des contextes admissibles}

À chaque couche \(\ell\), on maintient une banque de contextes :
\[
\mathcal B_\ell = \{\mathcal B_c^{(\ell)}\}_{c=1}^{C_\ell}.
\]

Chaque contexte \(c\) est associé à un hull convexe :
\[
\mathcal M_c^{(\ell)} = \operatorname{conv}\{\psi_\ell(v): v\in \mathcal B_c^{(\ell)}\},
\]
où \(\psi_\ell\) est une projection linéaire ou un petit module de compression de faible rang.

Le point important est que le contexte est géométrique, non attentionnel :
\[
c_i^{(\ell)}
=
\arg\min_{c\in\{1,\dots,C_\ell\}}
\operatorname{dist}\!\left(h_i^{(\ell)}, \mathcal M_c^{(\ell)}\right)^2.
\]

Le routeur est donc un routeur par distance de projection, cohérent avec le noyau PCH testé expérimentalement.

\subsection{Bloc PCH élémentaire}

Le bloc de couche \(\ell\) applique d'abord la purification contextuelle exacte :
\[
\tilde h_i^{(\ell)}
=
\Ppch_\ell\!\left(h_i^{(\ell)};c_i^{(\ell)}\right)
=
\arg\min_u
\left\{
\frac12\|u-h_i^{(\ell)}\|^2
+
\Phi_{\mathrm{PCH},c_i^{(\ell)}}^{(\ell)}(u)
\right\}.
\]

On prend
\[
\Phi_{\mathrm{PCH},c}^{(\ell)}(u)
=
\lambda_{\mathrm{ctx}}^{(\ell)}\operatorname{dist}(u,\mathcal M_c^{(\ell)})^2
+
\lambda_{\mathrm{ch}}^{(\ell)}\Omega_{\mathrm{ch}}^{(\ell)}(u)
+
\lambda_{\mathrm{reg}}^{(\ell)}\|u\|^2.
\]

Ensuite, un mélange local non attentionnel est appliqué :
\[
h_i^{(\ell+1)}
=
\tilde h_i^{(\ell)}
+
\mathrm{MLP}_\ell\!\left(\tilde h_i^{(\ell)}\right).
\]

La couche complète est donc :
\[
H^{(\ell)} \;\to\; C^{(\ell)} \;\to\; \tilde H^{(\ell)} \;\to\; H^{(\ell+1)}.
\]

\subsection{Interaction entre tokens sans attention}

L'architecture PCH pure n'utilise pas de pondérations pair-à-pair de type attention.  
L'échange d'information entre tokens passe par deux canaux seulement :

\begin{enumerate}[leftmargin=2em]
    \item un \emph{canal global} non attentionnel, par résumé :
    \[
    s^{(\ell)} = \frac1n\sum_{i=1}^n h_i^{(\ell)};
    \]
    \item un \emph{canal géométrique}, par attraction vers les hulls \(\mathcal M_c^{(\ell)}\).
\end{enumerate}

Le résumé global peut servir à sélectionner le contexte dominant, ou à moduler les paramètres de couche, mais il ne joue pas le rôle d'attention.

\subsection{Version diffusion : génération parallèle du texte}

Pour éviter un décodage autoregressif coûteux, on formule la génération comme une diffusion en espace latent.

Soit \(Z_0\in\mathbb R^{n\times d}\) le latent propre associé à la séquence cible.
Le processus direct est :
\[
Z_t
=
\sqrt{\bar\alpha_t}\,Z_0
+
\sqrt{1-\bar\alpha_t}\,\varepsilon,
\qquad
\varepsilon\sim\mathcal N(0,I).
\]

Le processus inverse est porté par un débruiteur latent \(\mu_\theta\), corrigé par PCH :
\[
Z_{t-1}
=
\arg\min_Z
\left\{
\frac1{2\sigma_t^2}\|Z-\mu_\theta(Z_t,t)\|_F^2
+
\Phi_{\mathrm{PCH},t}(Z)
\right\}.
\]

Ainsi, la génération n'est pas token-par-token au sens standard, mais \emph{séquentielle en temps de diffusion} et parallèle sur les positions.

\subsection{Objectifs d'entraînement}

L'entraînement est conçu en plusieurs termes complémentaires.

\paragraph{Perte de débruitage diffusion.}
\[
\mathcal L_{\mathrm{diff}}
=
\mathbb E_{x,t,\varepsilon}
\left[
\left\|
Z_0-\widehat Z_0(Z_t,t)
\right\|_F^2
\right].
\]

\paragraph{Perte de reconstruction textuelle.}
Pour les données dialogues, on ajoute une perte sur la réponse cible :
\[
\mathcal L_{\mathrm{tok}}
=
-\sum_{i\in \mathcal I_{\mathrm{resp}}}
\log p_\theta(x_i \mid H^{(L)}),
\]
où \(\mathcal I_{\mathrm{resp}}\) désigne les positions de la réponse à générer.

\paragraph{Perte PCH.}
On pénalise les sorties de couche par l'énergie exacte :
\[
\mathcal L_{\mathrm{PCH}}
=
\sum_{\ell=1}^L\sum_{i=1}^n
J_{\mathrm{PCH}}^{(\ell)}\!\left(
\tilde h_i^{(\ell)};
h_i^{(\ell)},
c_i^{(\ell)}
\right).
\]

\paragraph{Perte de routage géométrique.}
On peut ajouter un terme favorisant des choix de contexte stables :
\[
\mathcal L_{\mathrm{route}}
=
\sum_{\ell=1}^L\sum_{i=1}^n
\operatorname{dist}\!\left(h_i^{(\ell)},\mathcal M_{c_i^{(\ell)}}^{(\ell)}\right)^2.
\]

\paragraph{Perte totale.}
La perte finale est :
\[
\mathcal L
=
\mathcal L_{\mathrm{diff}}
+
\beta\,\mathcal L_{\mathrm{tok}}
+
\gamma\,\mathcal L_{\mathrm{PCH}}
+
\delta\,\mathcal L_{\mathrm{route}}.
\]

\subsection{Procédure d'entraînement complète}

L'entraînement se déroule en quatre phases.

\subsubsection{Phase 1 : initialisation}

On initialise :
\begin{itemize}[leftmargin=2em]
    \item l'embedding \(E_\theta\) ;
    \item les MLP de couche ;
    \item les projections \(\psi_\ell\) ;
    \item les banques de contextes \(\mathcal B_\ell\) ;
    \item le calendrier de diffusion \(\{\bar\alpha_t\}_{t=1}^T\).
\end{itemize}

Les banques de contextes peuvent être construites par pré-clustering sur un corpus de calibration, puis révisées périodiquement.

\subsubsection{Phase 2 : pré-calibration géométrique}

On extrait un ensemble de calibration par couche et par contexte, puis on forme les hulls :
\[
\mathcal M_c^{(\ell)}
=
\operatorname{conv}\{\psi_\ell(v):v\in\mathcal B_c^{(\ell)}\}.
\]

Cette phase fixe la géométrie admissible de départ.

\subsubsection{Phase 3 : entraînement conjoint}

Pour chaque mini-batch :
\begin{enumerate}[leftmargin=2em]
    \item tokeniser les séquences ;
    \item calculer les embeddings initiaux ;
    \item échantillonner un temps de diffusion \(t\) ;
    \item corrompre le latent en \(Z_t\) ;
    \item propager dans les blocs PCH de couches successives ;
    \item calculer la sortie de débruitage et la sortie textuelle ;
    \item calculer \(\mathcal L\) ;
    \item rétropropager à travers les solveurs convexes via différentiation implicite ou couche convexe différentiable ;
    \item mettre à jour les paramètres ;
    \item rafraîchir les banques de contextes à période fixée.
\end{enumerate}

\subsubsection{Phase 4 : stabilisation finale}

On fige partiellement ou totalement les banques de contextes, puis on affine les modules restants pour stabiliser :
\begin{itemize}[leftmargin=2em]
    \item la décroissance d'énergie ;
    \item la contraction locale ;
    \item la cohérence de génération ;
    \item la stabilité du routage.
\end{itemize}

\subsection{Différentiation à travers le prox exact}

Le prox PCH exact est un solveur convexe.  
Pendant l'entraînement, son gradient peut être obtenu par différentiation implicite du système KKT associé, ou par une couche convexe différentiable.

Si
\[
u^\star = \Ppch(w;c),
\]
alors les gradients par rapport aux paramètres appris sont calculés en utilisant la structure locale du problème convexe résolu exactement. Cela permet de garder le noyau PCH fidèle à sa définition tout en l'intégrant dans l'entraînement.

\subsection{Chargement du modèle entraîné}

Le checkpoint doit contenir au minimum :
\begin{itemize}[leftmargin=2em]
    \item \(E_\theta\) ;
    \item les MLP de couches ;
    \item les projections \(\psi_\ell\) ;
    \item les banques de contextes \(\mathcal B_\ell\) ;
    \item les paramètres \(\lambda_{\mathrm{ctx}}^{(\ell)},\lambda_{\mathrm{ch}}^{(\ell)},\lambda_{\mathrm{reg}}^{(\ell)}\) ;
    \item le calendrier de diffusion ;
    \item le tokenizer ;
    \item les statistiques de normalisation ;
    \item les éventuels tenseurs de décodage.
\end{itemize}

Le chargement consiste à restaurer exactement ces objets, puis à passer le modèle en mode inférence.

\subsection{Usage du modèle en chat prompt / réponse}

À l'usage, un échange chat est transformé en séquence structurée :
\[
(\texttt{system}, \texttt{user}_1, \texttt{assistant}_1, \dots, \texttt{user}_k).
\]

Le modèle reçoit :
\[
p = \text{historique complet + dernier message utilisateur},
\]
puis il génère la réponse \(r\).

La procédure d'inférence est :
\begin{enumerate}[leftmargin=2em]
    \item tokeniser le prompt \(p\) ;
    \item encoder les tokens ;
    \item construire l'état latent initial ;
    \item fixer le sous-espace de conditionnement du prompt ;
    \item initialiser la partie réponse par bruit latent ou par état neutre ;
    \item itérer la diffusion inverse PCH ;
    \item décoder la séquence réponse ;
    \item concaténer la réponse au dialogue courant ;
    \item répéter pour le tour suivant.
\end{enumerate}

Le point important est que le modèle ne consulte jamais l'historique via attention standard.  
Le contexte est porté par :
\begin{itemize}[leftmargin=2em]
    \item le prompt latent ;
    \item les banques de contextes ;
    \item les prox PCH ;
    \item la diffusion inverse.
\end{itemize}

\subsection{Pseudo-code de l'inférence chat}

\[
\boxed{
\begin{array}{l}
\text{Input : } p \\
\text{Tokenizer : } \tau(p) \\
\text{Embeddings : } H^{(0)} \\
\text{For }\ell=1\ldots L:\\
\qquad c_i^{(\ell)} \leftarrow \arg\min_c \operatorname{dist}(h_i^{(\ell)},\mathcal M_c^{(\ell)})^2\\
\qquad \tilde h_i^{(\ell)} \leftarrow \Ppch_\ell(h_i^{(\ell)};c_i^{(\ell)})\\
\qquad h_i^{(\ell+1)} \leftarrow \tilde h_i^{(\ell)}+\mathrm{MLP}_\ell(\tilde h_i^{(\ell)})\\
\text{Diffusion inverse : } Z_{T}\sim \mathcal N(0,I) \\
\text{For }t=T,\dots,1:\\
\qquad Z_{t-1}\leftarrow \arg\min_Z\left\{\frac1{2\sigma_t^2}\|Z-\mu_\theta(Z_t,t)\|_F^2+\Phi_{\mathrm{PCH},t}(Z)\right\} \\
\text{Decode : } r = \mathrm{Decoder}(Z_0) \\
\end{array}
}
\]

\subsection{Protocole de validation}

La validation d'un LLM PCH pur doit vérifier les points suivants :
\begin{enumerate}[leftmargin=2em]
    \item décroissance énergétique couche par couche ;
    \item résidus de point fixe faibles ;
    \item contraction locale non dégénérée ;
    \item stabilité du routage de contexte ;
    \item cohérence de la diffusion inverse ;
    \item qualité de génération sur tâches de chat ;
    \item absence totale d'attention standard dans le graphe computationnel.
\end{enumerate}

\subsection{Fidélité au noyau PCH}

L'architecture proposée reste fidèle au noyau PCH exact parce que :
\begin{itemize}[leftmargin=2em]
    \item chaque purification est un vrai minimiseur convexe ;
    \item le contexte est un hull convexe réel ;
    \item la contraction locale est mesurable ;
    \item la monotonie énergétique peut être contrôlée ;
    \item le routage est séparé du prox et ne le remplace pas.
\end{itemize}

\subsection{Résumé du pipeline complet}

Le pipeline final est :
\[
\text{texte}
\;\to\;
\text{tokenisation}
\;\to\;
E_\theta
\;\to\;
H^{(0)}
\;\to\;
\mathsf{PCHBlock}_1
\;\to\;
\cdots
\;\to\;
\mathsf{PCHBlock}_L
\;\to\;
\mathsf{DiffusionReverse}
\;\to\;
\text{décodeur}
\;\to\;
\text{réponse chat}.
\]

Aucune attention standard n'intervient à aucun point du pipeline.
Le modèle est donc un LLM à base de noyaux proximaux contextuels et de diffusion, entièrement compatible avec le formalisme PCH exact.

\end{document}